\section{R18: complete native seeded EMIT streams}
\label{sec:r18-journal}

R18 closes a concrete gap found on returning to the original seeded
workflow: a materialized artifact needs every fresh EMIT in order, including
those inside a fused chunk. The optional journal records that stream while
retaining the complete canonical transition, final receipts and ASA/NA/JK
feedback. Final-only paths remain separate.
\path{formal/SEEDED_EMIT_JOURNAL.md} specifies the complete contract;
the native schema is \code{ATOMOS-EMIT-JOURNAL-R1}.

\subsection{Transition and complete ordered stream}
For lane $i$, let $z_i(k)$ contain complete VM state, error, receipt and
optional feedback state at next dispatch epoch $k$. Preserve
\begin{equation}
 z_i(k+1)=C_k(z_i(k)),\qquad
 C_k=\operatorname{commit}_k\circ\operatorname{VM}_k
       \circ\operatorname{inject}_k.
\end{equation}
For pure execution, injection and word commit are identities. Append an
event exactly when the actual receipt says both executed and emitted.
A sticky state flag supplies no fresh event. EMIT-HALT supplies one;
identical payloads at distinct epochs remain distinct. A VM emission and
successful feedback commit are different observables, with final feedback
status recording a refused commit.

The native record is 40 bytes: a uint64 epoch and eight uint32 fields.
It records lane, pre-instruction cell, flags, payload, post-transition
lineage, cell, branch and status. Lineage, cell and status come from the
actual post-transition State64. This compact EMIT schema is distinct from
the vendor's full instruction trace.

With $E_i(k)$ denoting the record when present, define
\begin{equation}
 J_i[k,k+n)=\mathop{\Vert}_{j=0}^{n-1}
 \begin{cases}[E_i(k+j)]&\text{fresh VM emission},\\
 []&\text{otherwise},\end{cases}
 \label{eq:r18-stream}
\end{equation}
where $\Vert$ denotes ordered concatenation. At most one event per
transition gives $|J_i|\le n$ and proven total capacity $Ln$ for $L$ lanes.
Lane $i$ writes its disjoint interval $[in,(i+1)n)$ in increasing epoch
order. No global atomic append or GPU scheduling order defines output.
Host compaction publishes lane-major storage with offsets for each lane.

For contiguous unchanged execution,
\begin{equation}
 J_i[k,k+a+b)=J_i[k,k+a)\Vert J_i[k+a,k+a+b).
 \label{eq:r18-stream-composition}
\end{equation}
The R17 induction extends directly: equal pre-states give the same complete
transition, emission predicate, record, appended prefix and post-state.
Disjoint output intervals preserve this independently per lane. Across
several lanes, concatenate corresponding lane subsequences; appending whole
lane-major chunks gives a different global lane order.

\subsection{Admission, ownership and terminal status}
The VM and feedback journal APIs admit 0 through 256 requested epochs,
at least one lane, and representable epoch/allocation bounds. Zero work
observes an unchanged boundary without a new receipt or event.
\code{max\_events=0} selects automatic $Ln$ capacity; a nonzero value
below that bound refuses before transition or injection. A larger value
does not enlarge the actual allocation beyond $Ln$.

The event buffer initially uses $40Ln$ bytes on each of GPU and host,
plus lane counters, overflow indicators and complete initial/final boundary
arrays of size $O(L)$. Coupled execution captures both feedback states.
These allocations finish before corresponding mutable execution. The
immutable canonical source is retained separately. Runtime failure cannot
produce a successful journal, and no rollback is promised after it.

The synchronous drain checks capacity, lane, epoch bounds and ordering,
then captures all final VM words, errors and receipts, with applicable
feedback words. A nonzero VM error or feedback status gives
\code{Faulted}; otherwise all lanes halted gives \code{Complete}; otherwise
\code{Prefix}. A complete library journal may contain no emissions.
Artifact materialization separately requires at least one EMIT.
Every requested epoch runs even after early HALT/fault, preserving the
canonical final held receipt. Dispatch, execution, VM emission and accepted
feedback counts remain different.

Each chunk retains a shared immutable source object: process-unique owner
ID and every canonical \code{.tmg} byte read back from actual GPU-expanded
program words. It survives VM destruction. Move preserves the source;
reset changes generation. The owner ID is not a cryptographic identity.
Continuation requires the same source object, owner, generation, lane count,
execution mode, binding and complete word profile. It also requires
contiguous epochs and exact previous-final/next-initial VM states, errors
and applicable word states. Ordinary execution creates an explicit gap.
The feedback borrower additionally rejects external advance/reset, even
before a zero-work call.

Exclusive VM/lane ownership is required throughout the synchronous call,
with no concurrent same-lane work on another stream. Continuation checks
do not authenticate edited event rows or reconstruct a live owner from
serialized IDs. Pinned source and independent deterministic replay provide
the persisted output evidence.

\subsection{Literal byte materialization}
For flags $F$ and payload $W$, admit count
$c=(F\gg8)\mathbin{\&}7$ only for $1\le c\le4$. With endian bit
$b=(F\gg11)\mathbin{\&}1$, byte ordinal $j$, $0\le j<c$, is
\begin{equation}
 d_j=\begin{cases}j&b=0,\\c-1-j&b=1,\end{cases}
 \qquad B_j=(W\gg8d_j)\mathbin{\&}255.
\end{equation}
Big endian reverses the selected low $c$ bytes; it does not select the high
$c$ bytes. Concatenating decoded events in one lane's epoch order gives
the source-compatible artifact.

\code{wqk\_materialize} requires one canonical initial lane and source
bits 0 and 1, denoting seeded profile and emitted bytes. The horizon defaults
to the source header. Before mutable execution, it reserves all requested
events and up to four artifact bytes per epoch, and verifies GPU-expanded
source equality with the complete input. It joins contiguous chunks and
writes an artifact only after fault-free HALT with at least one EMIT.
Prefix/fault reports do not write an artifact. Source, artifact and report
paths must not alias.

The JSON contains all nine native event fields plus sequence, count and
byte order, with complete final State64 and run parameters. External replay
checks all twelve fields, lane zero, complete state and output bytes, and
requires literal true artifact-written/source-equality assertions. These
assertions are not authenticated credentials; pinned bytes and exact replay
supply the substantive output comparison.

CLI timings describe host setup; journal allocation, execution, transfers,
compaction and decoding; artifact write; and total through artifact
persistence. JSON encoding/report output is excluded. They establish no
device-only or matched external materializer speedup.

The original compiler still validates dependencies, evaluates any formal
tree on the host, and lowers the result into Cell48 EMIT cells. R18 executes
and materializes those cells natively. GPU \code{formal.evaluate}, arbitrary
JSON evaluation and native learner search remain separate work. This exact
stream contract adds no physical law or continuous precision.
\clearpage
